\chapter{Конструкторская часть}

В данном разделе представлены требования к разрабатываемому программному обеспечению для трехмерного моделирования движения автомобиля, а также детально описаны алгоритмы и структуры данных, выбранные для реализации системы визуализации.

\section{Требования к программному обеспечению}

Разрабатываемая программа должна предоставлять пользователю графический интерфейс со следующим функционалом:
\begin{itemize}
	\item загрузка и отображение трехмерной сцены городской среды из конфигурационных файлов;
	\item управление положением и ориентацией виртуальной камеры в пространстве;
	\item запуск анимации движения автомобиля по автоматически построенному маршруту;
	\item отображение текущих параметров производительности (кадров в секунду).
\end{itemize}

Разработанное ПО должно соответствовать следующим техническим требованиям:
\begin{itemize}
	\item использование алгоритма Z-буфера для корректного удаления невидимых поверхностей;
	\item реализация метода карт теней (Shadow Mapping) для построения динамических теней;
	\item применение модели освещения Ламберта с плоским закрашиванием (Flat Shading);
	\item поддержка текстурирования объектов сцены;
	\item расчет маршрута движения с учетом графа дорожной сети и правил дорожного движения.
\end{itemize}

\section{Описание структур данных}

Для реализации алгоритмов визуализации и моделирования движения используются следующие ключевые структуры данных.

\begin{enumerate}
	\item \textbf{Сцена (Scene)} — контейнер верхнего уровня, содержащий списки всех статических (здания, дороги) и динамических (автомобиль) объектов, параметры источника света и объект камеры.
	
	\item \textbf{Геометрический объект (SceneObject)} включает в себя:
	\begin{itemize}
		\item массив вершин (\texttt{dots}) — список точек в трехмерном пространстве $(x, y, z)$;
		\item массив граней (\texttt{faces}) — список полигонов, где каждый полигон задается индексами вершин, цветом и текстурными координатами;
		\item вектор положения центра объекта в мировых координатах;
		\item матрицу поворота, определяющую ориентацию объекта;
		\item вектор направления «вперед» (\texttt{forward}) для корректного движения.
	\end{itemize}
	
	\item \textbf{Камера (Camera)} содержит:
	\begin{itemize}
		\item вектор положения в мировом пространстве;
		\item вектор направления взгляда (\texttt{target});
		\item вектор ориентации «вверх» (\texttt{up});
		\item матрицы вида (\texttt{view\_matrix}) и проекции, пересчитываемые при изменении параметров.
	\end{itemize}
	
	\item \textbf{Источник света (ShadowCaster)} включает:
	\begin{itemize}
		\item вектор направления света (для направленного источника);
		\item параметры интенсивности и цвета;
		\item собственные матрицы вида и проекции для генерации карты теней;
		\item буфер глубины (Shadow Map) — двумерный массив для хранения расстояний от источника света до поверхностей.
	\end{itemize}
	
	\item \textbf{Дорожная сеть} представлена двумя матрицами:
	\begin{itemize}
		\item матрица проходимости (бинарная: 0 — препятствие, 1 — дорога);
		\item матрица направлений (битовая маска разрешенных направлений движения из каждой клетки).
	\end{itemize}
\end{enumerate}

\section{Общий алгоритм построения изображения}

Процесс генерации кадра состоит из двух основных проходов: прохода для генерации карты теней и прохода для финальной отрисовки сцены с точки зрения камеры.

На вход алгоритма подаются списки объектов сцены и параметры камеры. На выходе формируется двумерное растровое изображение, которое выводится на экран.

\textbf{Этап 1: Генерация карты теней (Shadow Pass)}
\begin{enumerate}
	\item Устанавливается камера в положение источника света.
	\item Используется ортогональная проекция для имитации солнечных лучей.
	\item Производится растеризация всех объектов сцены, но вместо цвета в буфер (Shadow Map) записывается только глубина ($z$-координата) ближайшего фрагмента.
\end{enumerate}

\textbf{Этап 2: Финальная отрисовка (Camera Pass)}
\begin{enumerate}
	\item Устанавливается камера наблюдателя.
	\item Инициализируются буфер цвета и Z-буфер (заполняется значением бесконечности).
	\item Для каждого объекта сцены выполняется конвейер визуализации:
	\begin{itemize}
		\item трансформация вершин (Model-View-Projection);
		\item отсечение невидимых полигонов (отключено для корректной работы с разнородной геометрией);
		\item растеризация полигонов методом построчного сканирования (Scanline);
		\item для каждого пикселя выполняется тест глубины (Z-test), расчет освещения и наложение текстур.
	\end{itemize}
\end{enumerate}

\section{Аффинные преобразования}

Для размещения объектов в сцене и анимации движения используются матрицы аффинных преобразований размером $4 \times 4$ в системе однородных координат.

Матрица поворота вокруг оси $Z$ на угол $\alpha$ (используется для поворота автомобиля):
\begin{equation}
	R_z(\alpha) = \begin{pmatrix}
		\cos \alpha & -\sin \alpha & 0 & 0 \\
		\sin \alpha & \cos \alpha & 0 & 0 \\
		0 & 0 & 1 & 0 \\
		0 & 0 & 0 & 1
	\end{pmatrix}
\end{equation}

Матрица переноса на вектор $(dx, dy, dz)$:
\begin{equation}
	T(dx, dy, dz) = \begin{pmatrix}
		1 & 0 & 0 & 0 \\
		0 & 1 & 0 & 0 \\
		0 & 0 & 1 & 0 \\
		dx & dy & dz & 1
	\end{pmatrix}
\end{equation}

Итоговое положение вершины $V_{new}$ вычисляется как произведение исходной вершины $V$ на композицию матриц трансформации. В программной реализации для ускорения вычислений используется библиотека NumPy, позволяющая выполнять операции над массивами вершин векторизованно: $V_{new} = V_{array} \cdot M$.

\section{Приведение к пространству камеры}

Для отображения сцены с точки зрения наблюдателя все объекты переводятся в систему координат камеры. Матрица вида ($View Matrix$) формируется на основе позиции камеры $P$, точки направления взгляда $T$ и вектора «вверх» $U$.

Базис камеры рассчитывается следующим образом:
\begin{equation}
	D = \frac{T - P}{||T - P||}, \quad R = \frac{D \times U}{||D \times U||}, \quad U' = R \times D
\end{equation}
где $D$ — вектор направления, $R$ — вектор «вправо», $U'$ — скорректированный вектор «вверх».

Матрица вида представляет собой комбинацию поворота и переноса, переводящую мир в систему отсчета, где камера находится в начале координат и направлена вдоль оси $Z$.

\section{Перспективная проекция}

Для создания эффекта перспективы используется матрица перспективной проекции. Она преобразует усеченную пирамиду видимости в куб в нормализованных координатах устройства.

Ключевым этапом является \textbf{перспективное деление}: координаты $x$ и $y$ делятся на глубину $z$.
\begin{equation}
	x_{screen} = x \cdot \frac{d}{z + d} \cdot \text{scale} + x_{offset}
\end{equation}
\begin{equation}
	y_{screen} = y \cdot \frac{d}{z + d} \cdot \text{scale} + y_{offset}
\end{equation}
где $d$ — фокусное расстояние, определяющее степень перспективного искажения.

\section{Алгоритм построчного сканирования с Z-буфером}

Для растеризации полигонов (преобразования векторных треугольников в пиксели) используется алгоритм построчного сканирования (Scanline), совмещенный с Z-буфером для удаления невидимых поверхностей.

\textbf{Схема работы алгоритма:}
\begin{enumerate}
	\item Вершины треугольника проецируются на экранную плоскость.
	\item Определяется ограничивающий прямоугольник полигона по оси Y.
	\item Происходит итерация по строкам экрана ($y$) от минимальной до максимальной координаты $Y$ полигона.
	\item Для каждой строки находятся точки пересечения сканирующей линии с ребрами треугольника ($x_{start}, x_{end}$).
	\item Внутри интервала $[x_{start}, x_{end}]$ для каждого пикселя интерполируется глубина $z$.
	\item \textbf{Z-тест:} Значение $z$ сравнивается с текущим значением в буфере глубины $z_{buffer}[y][x]$.
	\begin{itemize}
		\item Если $z < z_{buffer}[y][x]$, то пиксель считается видимым. В буфер записывается новое значение $z$, а в буфер кадра — вычисленный цвет пикселя.
		\item Иначе пиксель отбрасывается.
	\end{itemize}
\end{enumerate}

Благодаря использованию библиотеки NumPy, внутренний цикл по $X$ выполняется векторизованно, обрабатывая целую строку пикселей одной операцией, что обеспечивает высокую производительность на интерпретируемом языке Python.

\section{Расчет освещения и теней}

Для закраски пикселей используется модель плоского закрашивания (Flat Shading) с учетом освещенности по закону Ламберта и теней.

\begin{enumerate}
	\item \textbf{Нормаль поверхности:} Для каждого треугольника вычисляется вектор нормали $N$ как векторное произведение двух его сторон.
	
	\item \textbf{Диффузное освещение:} Коэффициент освещенности $I_{diffuse}$ рассчитывается как скалярное произведение нормали $N$ и вектора направления на свет $L$:
	\begin{equation}
		I_{diffuse} = |N \cdot L|
	\end{equation}
	Использование модуля позволяет освещать грани с обеих сторон (актуально для стен без толщины).
	
	\item \textbf{Расчет тени (Shadow Mapping):}
	\begin{itemize}
		\item Координаты текущего пикселя $(x, y, z)$ переводятся в пространство источника света.
		\item Полученные координаты $(u_{light}, v_{light})$ используются для выборки глубины из карты теней: $d_{shadow} = \text{ShadowMap}[v_{light}, u_{light}]$.
		\item Текущая глубина пикселя в пространстве света $z_{light}$ сравнивается с $d_{shadow}$.
		\item Если $z_{light} > d_{shadow} + \text{bias}$, то точка находится в тени. Коэффициент освещенности уменьшается до значения фонового света.
	\end{itemize}
\end{enumerate}

Итоговый цвет пикселя вычисляется как произведение цвета текстуры на результирующий коэффициент освещенности.

\section{Организация поиска пути}

Для моделирования движения автомобиля используется \textbf{волновой алгоритм} на ориентированном графе.

\begin{enumerate}
	\item Дорожная карта дискретизируется в сетку.
	\item Каждой ячейке сопоставляется битовая маска разрешенных направлений (Вверх, Вниз, Влево, Вправо).
	\item Алгоритм начинает распространение волны от стартовой точки. При переходе в соседнюю клетку проверяется битовая маска. Это гарантирует соблюдение правил дорожного движения (например, запрет движения против потока).
	\item После достижения целевой точки путь восстанавливается обратным ходом от финиша к старту.
\end{enumerate}

Найденный путь преобразуется в список мировых координат, между которыми происходит плавная интерполяция положения и угла поворота автомобиля при анимации.

\section{Вывод}

В разделе разработаны требования к ПО, определены ключевые структуры данных для хранения сцены. Детально описан конвейер визуализации, включающий матричные преобразования, алгоритм построчного сканирования с Z-буфером, метод карт теней и модель освещения Ламберта. Также представлена логика организации движения на основе волнового алгоритма. Выбранные решения оптимизированы для программной реализации и обеспечивают корректное отображение трехмерной сцены в реальном времени.